% Bài làm buổi 20/10 — Cam
\begin{problem}
Phân tích các đa thức sau thành nhân tử: a) $x^2 - xy + 3x - 3y$; b) $2xy + 3z + 6y + xz$.
a) $x^2 - xy + 3x - 3y = (x^2 - xy) + (3x - 3y) = x(x - y) + 3(x - y) = (x - y)(x + 3)$
b) $2xy + 3z + 6y + xz = (2xy + 6y) + (xz + 3z) = 2y(x + 3) + z(x + 3) = (x + 3)(2y + z)$
\end{problem}

\begin{problem}[8DS-PS-003]
ĐKXĐ: $x \neq y$
$\dfrac{x^2}{x - y} + \dfrac{y^2}{y - x} = \dfrac{x^2}{x - y} + \dfrac{y^2}{x - y}$
$= \dfrac{x^2 + y^2}{x - y}$
Vậy kết quả là $\dfrac{x^2 + y^2}{x - y}$.
\end{problem}

\begin{problem}[8DS-PT1-002]
a) Thay $x = -2$: vế trái $3 \cdot (-2) + 5 = -1$, vế phải $-2 + 1 = -1$. Vậy $x = -2$ là nghiệm.
Thay $x = 1$: vế trái $3 \cdot 1 + 5 = 8$, vế phải $1 + 1 = 2$. Vì $8 \neq 2$ nên $x = 1$ không là nghiệm.
b) Thay $x = 3$ vào phương trình: $2 \cdot 3 + m = 5 - 3$
$6 + m = 2$
$m = -4$
Vậy $m = -4$.
\end{problem}

\begin{problem}[8DS-HSBN-003]
a) Thay $x = 1$: $y = -2 \cdot 1 + 3 = 1$ nên $M(1; 1) \subset d$.
Thay $x = 2$: $y = -2 \cdot 2 + 3 = -1 \neq 1$ nên $N(2; 1) \notin d$.
b) Cho $x = 0$ thì $y = 3$, vậy $A(0; 3)$.
Cho $y = 0$ thì $-2x + 3 = 0$, $x = \dfrac{3}{2}$, vậy $B\left(\dfrac{3}{2}; 0\right)$.
Đường thẳng $d$ là đường thẳng đi qua hai điểm $A$ và $B$.
c) $S_{OAB} = \dfrac{1}{2} \cdot OA \cdot OB = \dfrac{1}{2} \cdot 3 \cdot \dfrac{3}{2} = \dfrac{9}{4}$
\end{problem}
